\section{模型的评价、改进与推广}

\subsection{模型的优点}

\begin{enumerate}
\item \textbf{广义标度律结构简洁且物理自洽。}以“可约损失 × 配比因子”的乘性结构统一了参数、数据、质量与配比四个因素：$Q=1$、参考配比时严格退化为经典标度律，任何资源组合下都有 $L>E$，并由此严格证明了配比与规模—质量决策的可分性（定理~\ref{thm:gsl}）。配比效应尺度不变这一关键假设不是凭空设定的，而是在附件 A 的三个数量级规模上得到了检验（$\theta$ 变异系数 $0.17$）。
\item \textbf{关键结论有闭式解与严格证明。}质量—规模替代的充要条件、等价增幅与临界规模（定理~\ref{thm:subst}），上下文临界长度 $6/\eta$，计算最优指数 $N^*\propto C^{0.452}$，均为解析结果；数值结果与闭式解相互印证。
\item \textbf{结构性转移有严格的数学定义。}以 KKT 活跃约束集合的变化定义主要转移、以配置弹性断点定义次要转移，并以参数抽样下的出现概率界定稳定性，使“是否发生质变”成为可以计算、可以检验的问题。
\item \textbf{验证体系完整，数据使用规范。}每个模型都在未参与训练的数据上检验：配比模型用三个规模的独立检验集，经典标度律用留一规模与跨族数据，质量项用四种方式的成块留出，桥接模型用来源族留一，前沿模型用滚动起点回测。半合成与估算数据只用于补充分析与边界检验。
\item \textbf{不确定性全链条传播。}从问题一的配比 bootstrap、问题二的参数 bootstrap，到问题三的最优解区间、问题四的桥接与前沿联合抽样，每个关键数值都附有区间，并识别出能力结论的主导误差来源（Loss--Benchmark 映射）。
\end{enumerate}

\subsection{模型的不足}

\begin{enumerate}
\item \textbf{推断领域的质量无法细分。}质量数据只覆盖 7 个领域，11 个配比领域的质量只能取公共参考值；配比效应本身由实验数据直接估计，因此不影响配比结论，但限制了“质量 × 配比”交互的刻画。
\item \textbf{跨体系对齐依赖锚定假设。}问题一与附件 B 的质量尺度以 $Q_{\mathrm{ref}}$ 为锚点线性对齐，$s_Q$ 无法由数据唯一识别；敏感性分析显示 $s_Q$ 是影响高预算最优 Loss 的首要因素。
\item \textbf{高预算结论依赖外推。}$10^{24}$ FLOPs 的最优规模远超 B1/B7 的观测范围，严格支持域下 Loss 会上升约 $0.18$；质量成本函数为题设情景，而非实测工程成本。
\item \textbf{历史面板较短。}确认性前沿只有 7 个月、41 个模型，限制了非规模技术状态的识别能力；Loss--Benchmark 桥接的留族误差约 10 分，能力的绝对水平精度有限。
\end{enumerate}

\subsection{模型的改进}

\begin{enumerate}
\item 在同一模型族与验证集上设计 $(N,D,Q,\boldsymbol p)$ 的分层因子实验，直接识别质量映射斜率 $s_Q$ 与配比系数 $\theta$，并检验质量与配比之间是否存在交互。
\item 收集语料过滤、去重、合成与审核的实际成本，替代题设质量成本函数；把显存、通信与硬件利用率纳入约束，把单一 FLOPs 约束扩展为多资源约束，并考虑推理成本对最优规模的影响\cite{sardana2024beyond}。
\item 对超出支持域的配置采用分布鲁棒优化，在标度律参数的不确定集上优化最坏情况 Loss，以降低外推风险。
\item 延长同口径能力面板，加入架构、数据处理、后训练算力等可观测特征，把“非规模技术状态”拆分为可解释的技术通道；引入观测标度律\cite{ruan2024observational}的低维能力因子改进 Loss--Benchmark 桥接。
\end{enumerate}

\subsection{模型的推广}

本文“数据度量—广义标度律—约束优化—能力预测”的框架不依赖特定模型。对视觉、多模态或科学计算基础模型，只需把 Token 替换为图像块、样本或模态单元，重新定义质量指标与序列相关开销，即可沿用同一建模路径。广义标度律的乘性配比结构与可分性定理，可直接用于数据配比实验的设计：先在小模型上确定配比，再把配比迁移到大模型，只需关注规模与质量的分配。在训练管理中，该框架可作为算力预算的规划工具：以中心解给出基准配置，以结构性转移点判断何时应启动质量投入，以外推标签提示决策风险。
